% questions.tex —— 2014 年·全国大纲卷·理科（广西） · 题目内容片段（无答案）
% 由 tools/generate_projects.py 依来源总账自 2014/2014_数学真题_大纲理_试题.md 生成，勿手改（改源 MD 后重跑）


\begin{gknotice}
\textbf{来源说明}：本题面依据固定源 PDF 的两页 A3 题面逐题转录，原卷印刷页码为 869—870。第 869 面为选择题 1—12、填空题 13—16、解答题 17—19；第 870 面为解答题 20—22。

\textbf{图形说明}：第 19 题三棱柱图由项目内 TikZ 矢量源重绘；其余各题原页无图。

\textbf{解析归属}：固定源不附答案或解析；配套答案与解析均由本库按本题面独立复算。
\end{gknotice}

\gkpart{一、选择题}

\begin{question}
\gkti{1}设 $z=\dfrac{10\mathrm i}{3+\mathrm i}$，则 $z$ 的共轭复数为\mbox{（　　）}

（A）$-1+3\mathrm i$；（B）$-1-3\mathrm i$；（C）$1+3\mathrm i$；（D）$1-3\mathrm i$。

\end{question}
\begin{question}
\gkti{2}设集合 $M=\{x\mid x^2-3x-4<0\}$，$N=\{x\mid 0\leqslant x\leqslant 5\}$，则 $M\cap N=$\mbox{（　　）}

（A）$(0,4]$；（B）$[0,4)$；（C）$[-1,0)$；（D）$(-1,0]$。

\end{question}
\begin{question}
\gkti{3}设 $a=\sin33^\circ$，$b=\cos55^\circ$，$c=\tan35^\circ$，则\mbox{（　　）}

（A）$a>b>c$；（B）$b>c>a$；（C）$c>b>a$；（D）$c>a>b$。

\end{question}
\begin{question}
\gkti{4}若向量 $\boldsymbol a$、$\boldsymbol b$ 满足：$|\boldsymbol a|=1$，$(\boldsymbol a+\boldsymbol b)\perp\boldsymbol a$，$(2\boldsymbol a+\boldsymbol b)\perp\boldsymbol b$，则 $|\boldsymbol b|=$\mbox{（　　）}

（A）$2$；（B）$\sqrt2$；（C）$1$；（D）$\dfrac{\sqrt2}{2}$。

\end{question}
\begin{question}
\gkti{5}有 6 名男医生、5 名女医生，从中选出 2 名男医生、1 名女医生组成一个医疗小组，则不同的选法共有\mbox{（　　）}

（A）60 种；（B）70 种；（C）75 种；（D）150 种。

\end{question}
\begin{question}
\gkti{6}已知椭圆 $C:\dfrac{x^2}{a^2}+\dfrac{y^2}{b^2}=1$（$a>b>0$）的左、右焦点分别为 $F_1$、$F_2$，离心率为 $\dfrac{\sqrt3}{3}$，过 $F_2$ 的直线 $l$ 交 $C$ 于 $A$、$B$ 两点，若 $\triangle AF_1B$ 的周长为 $4\sqrt3$，则 $C$ 的方程为\mbox{（　　）}

（A）$\dfrac{x^2}{3}+\dfrac{y^2}{2}=1$；（B）$\dfrac{x^2}{3}+y^2=1$；（C）$\dfrac{x^2}{12}+\dfrac{y^2}{8}=1$；（D）$\dfrac{x^2}{12}+\dfrac{y^2}{4}=1$。

\end{question}
\begin{question}
\gkti{7}曲线 $y=x\mathrm e^{x-1}$ 在点 $(1,1)$ 处切线的斜率等于\mbox{（　　）}

（A）$2\mathrm e$；（B）$\mathrm e$；（C）$2$；（D）$1$。

\end{question}
\begin{question}
\gkti{8}正四棱锥的顶点都在同一球面上，若该棱锥的高为 $4$，底面边长为 $2$，则该球的表面积为\mbox{（　　）}

（A）$\dfrac{81\pi}{4}$；（B）$16\pi$；（C）$9\pi$；（D）$\dfrac{27\pi}{4}$。

\end{question}
\begin{question}
\gkti{9}已知双曲线 $C$ 的离心率为 $2$，焦点为 $F_1$、$F_2$，点 $A$ 在 $C$ 上，若 $|F_1A|=2|F_2A|$，则 $\cos\angle AF_2F_1=$\mbox{（　　）}

（A）$\dfrac14$；（B）$\dfrac13$；（C）$\dfrac{\sqrt2}{4}$；（D）$\dfrac{\sqrt2}{3}$。

\end{question}
\begin{question}
\gkti{10}等比数列 $\{a_n\}$ 中，$a_4=2$，$a_5=5$，则数列 $\{\lg a_n\}$ 的前 $8$ 项和等于\mbox{（　　）}

（A）$6$；（B）$5$；（C）$4$；（D）$3$。

\end{question}
\begin{question}
\gkti{11}已知二面角 $\alpha-l-\beta$ 为 $60^\circ$，$AB\subset\alpha$，$AB\perp l$，$A$ 为垂足，$CD\subset\beta$，$C\in l$，$\angle ACD=135^\circ$，则异面直线 $AB$ 与 $CD$ 所成角的余弦值为\mbox{（　　）}

（A）$\dfrac14$；（B）$\dfrac{\sqrt2}{4}$；（C）$\dfrac{\sqrt3}{4}$；（D）$\dfrac12$。

\end{question}
\begin{question}
\gkti{12}函数 $y=f(x)$ 的图象与函数 $y=g(x)$ 的图象关于直线 $x+y=0$ 对称，则 $y=f(x)$ 的反函数是\mbox{（　　）}

（A）$y=g(x)$；（B）$y=g(-x)$；（C）$y=-g(x)$；（D）$y=-g(-x)$。

\end{question}
\gkpart{二、填空题}

\begin{question}
\gkti{13}$\left(\dfrac{x}{\sqrt y}-\dfrac{y}{\sqrt x}\right)^8$ 的展开式中 $x^2y^2$ 的系数为\gkblank{}。（用数字作答）

\end{question}
\begin{question}
\gkti{14}设 $x$、$y$ 满足约束条件
\nopagebreak[4]
\[\begin{cases}
x-y\geqslant 0,\\
x+2y\leqslant 3,\\
x-2y\leqslant 1,
\end{cases}\]
则 $z=x+4y$ 的最大值为\gkblank{}。

\end{question}
\begin{question}
\gkti{15}直线 $l_1$ 和 $l_2$ 是圆 $x^2+y^2=2$ 的两条切线，若 $l_1$ 与 $l_2$ 的交点为 $(1,3)$，则 $l_1$ 与 $l_2$ 的夹角的正切值等于\gkblank{}。

\end{question}
\begin{question}
\gkti{16}若函数 $f(x)=\cos2x+a\sin x$ 在区间 $\left(\dfrac\pi6,\dfrac\pi2\right)$ 上是减函数，则 $a$ 的取值范围是\gkblank{}。

\end{question}
\gkpart{三、解答题}

\begin{question}
\gkti{17}$\triangle ABC$ 的内角 $A$、$B$、$C$ 的对边分别为 $a$、$b$、$c$，已知 $3a\cos C=2c\cos A$，$\tan A=\dfrac13$，求 $B$。

\end{question}
\begin{question}
\gkti{18}等差数列 $\{a_n\}$ 的前 $n$ 项和为 $S_n$，已知 $a_1=10$，$a_2$ 为整数，且 $S_n\leqslant S_4$。

（1） 求 $\{a_n\}$ 的通项公式；

（2） 设 $b_n=\dfrac{1}{a_na_{n+1}}$，求数列 $\{b_n\}$ 的前 $n$ 项和 $T_n$。

\Needspace{0.4\textheight}
\end{question}
\begin{question}
\gkti{19}如图，三棱柱 $ABC-A_1B_1C_1$ 中，点 $A_1$ 在平面 $ABC$ 内的射影 $D$ 在 $AC$ 上，$\angle ACB=90^\circ$，$BC=1$，$AC=CC_1=2$。
\nopagebreak[4]
\begin{center}% ============================================================
% f2014_outline_science_q19_prism.tex
% 2014 全国大纲卷（理科）第 19 题：三棱柱 ABC-A1B1C1
% A1 在平面 ABC 内的射影 D 在 AC 上，∠ACB=90°。
% 来源 PDF：普通高考/2014/2014大纲理(广西).pdf
% 原页：第 1 页（A3 横排）第 3 栏，印面 869，第 19 题题干下方
% 画法为平行投影示意，不按比例；虚线为不可见棱及辅助线。
% ============================================================
\begin{tikzpicture}[>=Stealth, line width=0.7pt, font=\small,
  line cap=round, line join=round]
  % 底面 ABC：C 为直角顶点，CA=2、CB=1；D 为 AC 中点，
  % A1 在 D 正上方（h=√3），三棱柱侧棱 AA1∥CC1∥BB1。
  \coordinate (A)  at (0.00,0.00);
  \coordinate (C)  at (1.80,1.15);
  \coordinate (B)  at (3.90,1.00);
  \coordinate (D)  at (0.90,0.58);
  \coordinate (A1) at (0.95,3.38);
  \coordinate (C1) at (2.75,4.53);
  \coordinate (B1) at (4.85,4.38);

  % 可见棱
  \draw (A) -- (B) -- (B1) -- (C1) -- (A1) -- cycle;
  \draw (A) -- (A1);
  \draw (A1) -- (B1);
  \draw (A1) -- (B);

  % 不可见棱与辅助线
  \draw[dashed] (A) -- (C) -- (B);
  \draw[dashed] (C) -- (C1);
  \draw[dashed] (A1) -- (D);
  \draw[dashed] (A1) -- (C);

  \fill (A) circle (1.1pt) node[below left] {$A$};
  \fill (B) circle (1.1pt) node[right] {$B$};
  \fill (C) circle (1.1pt) node[above left] {$C$};
  \fill (D) circle (1.1pt) node[below=1pt] {$D$};
  \fill (A1) circle (1.1pt) node[left] {$A_1$};
  \fill (C1) circle (1.1pt) node[above] {$C_1$};
  \fill (B1) circle (1.1pt) node[above right] {$B_1$};
\end{tikzpicture}
\end{center}
\nopagebreak[4]
（1） 证明：$AC_1\perp A_1B$；

（2） 设直线 $AA_1$ 与平面 $BCC_1B_1$ 的距离为 $\sqrt3$，求二面角 $A_1-AB-C$ 的大小。

\end{question}
\begin{question}
\gkti{20}设每个工作日甲、乙、丙、丁 4 人需使用某种设备的概率分别为 $0.6$、$0.5$、$0.5$、$0.4$，各人是否需使用设备相互独立。

（1） 求同一工作日至少 3 人需使用设备的概率；

（2） $X$ 表示同一工作日需使用设备的人数，求 $X$ 的数学期望。

\end{question}
\begin{question}
\gkti{21}已知抛物线 $C:y^2=2px$（$p>0$）的焦点为 $F$，直线 $y=4$ 与 $y$ 轴的交点为 $P$，与 $C$ 的交点为 $Q$，且 $|QF|=\dfrac54|PQ|$。

（1） 求 $C$ 的方程；

（2） 过 $F$ 的直线 $l$ 与 $C$ 相交于 $A$、$B$ 两点，若 $AB$ 的垂直平分线 $l'$ 与 $C$ 相交于 $M$、$N$ 两点，且 $A$、$M$、$B$、$N$ 四点在同一圆上，求 $l$ 的方程。

\end{question}
\begin{question}
\gkti{22}函数 $f(x)=\ln(x+1)-\dfrac{ax}{x+a}$（$a>1$）。

（1） 讨论 $f(x)$ 的单调性；

（2） 设 $a_1=1$，$a_{n+1}=\ln(a_n+1)$，证明：$\dfrac{2}{n+2}<a_n\leqslant\dfrac{3}{n+2}$。

\end{question}